\section{Introduction}
\label{sec:introduction}

Language-model agents use tools to inspect schemas, execute Python or SQL, retrieve documentation, and produce analytical reports \citep{schick2023toolformer,yao2023react,qin2023toolllm}. Their answers depend on both the computation and how they interpret its output. Successful execution alone does not establish that the returned answer is correct.

Data workflows encounter errors that preserve a normal-looking interface: an aggregate uses the wrong denominator, a statistic is attached to the wrong entity, or a retrieved data dictionary is stale. These errors can pass format checks and still change the answer. Recent benchmarks study silent tool errors and recovery under unreliable tool environments \citep{sun2024toolsfail,tian2026toolbenchx}. We focus on how data agents use a corrupted observation when the source data remain available for checking. We call this setting \emph{silent tool poisoning}: the tool executes successfully, but the observation supplied to the agent contains controlled factual or semantic corruption.

ToxicBench asks whether agents detect, verify, or adopt this corrupted evidence. Each task pairs the same query, data, model, and adapter in clean and poisoned environments. Numerical operators alter magnitudes, signs, ratios, or rankings; semantic operators alter entity bindings, column meanings, metadata, or retrieved evidence. Alongside task success, trajectory metrics separate adoption without checking from adoption after verification and successful recovery. This distinction matters because an agent may request another computation yet still select the wrong conclusion.

High clean-task success does not ensure robustness to these changes. Across LangGraph, smolagents, and AutoGen in the 118-task GPT evaluation, poisoned task success falls by 26--39 percentage points. Cross-model evaluations extend the comparison to Claude and Qwen. The paired design makes it possible to examine this degradation alongside the agent's checking behavior and final answer, rather than treating every failure as a lack of task-solving ability.

Our verification study separates the value of another attempt from an explicit instruction to check. Double-pass, Verification-only, and Generic Guard share route counts, per-route step limits, and per-request output-token caps. Under one-shot poisoning, an ordinary second attempt can access later uncorrupted observations; under repeated poisoning, fresh checks may also be corrupted. Controlled experiments on three public tables test the effect of these later evidence conditions. A post-freeze human evaluation corroborates the retry benefit over Base on the audited tasks and confirms that checking can coexist with adoption of a poisoned answer.

Our contributions are:
\begin{itemize}
    \item An inspectable data-analysis benchmark that changes returned evidence while preserving source tables, including numerical, semantic/schema, and multi-table tasks.
    \item A distinction between checking and answer adoption: trace-level and human evidence show that a fresh check need not prevent adoption of a poisoned conclusion, especially when the check is also corrupted.
    \item Controls separating another attempt, verification instructions, and later evidence conditions, including field-bound, matched-parent comparisons on three public tables that identify conditions supporting recovery.
\end{itemize}
